\subsection{Robustness and the Background Budget}\label{sec:res-robust}

\paragraph{Attacker seeds and data splits}
Ground truth is itself an estimate: attackers are trained with random initializations and
early-stopping splits, and $\Vy$ is measured on one random test split. We repeated the complete
three-attacker retraining with two further seeds (RTS-GMLC, three targets; PJM load) and on two
further random 70/30 splits (RTS-GMLC). Across these reruns the mean absolute change in $\Vy$ is
0.001--0.008, the Spearman correlation of $M^{(2)}$ across fields with the main run is 0.983--0.999,
the set of $\tau$-critical fields is identical at $\tau=0.5$ (Jaccard 1.0 in every case) and agrees
with Jaccard 0.93--1.00 at $\tau=0.7$, and the number of combination-critical fields changes by at
most one or two fields (Fig.~\ref{fig:robust}c). The conclusions of
Sections~\ref{sec:res-amp}--\ref{sec:res-def} are therefore not artefacts of a particular seed or
split.

\paragraph{Threshold}
Fig.~\ref{fig:robust}(d) sweeps $\tau$ from 0.3 to 0.9 over all ten targets. At every threshold the
share of fields that are safe alone but critical with at most two background fields exceeds the share
of fields that are unsafe alone (e.g.\ 43\% versus 10\% at $\tau=0.7$). The gap between the two grading
philosophies is thus not tied to the choice of $\tau$, although its size is
(Section~\ref{sec:application}).

\paragraph{Choice of $K$}
Exhaustive retraining of all sets of four fields on every data set (for NEM an additional 27,405 sets)
gives the exact $M^{(3)}$. Measured by the \emph{marginal} risk, $K=2$ is close to saturation:
$M^{(2)}/M^{(3)}$, the ratio of mean field risk, lies between 0.958 and 1.000 for nine targets and is
0.813 for NEM TUMUT3, a hydro unit whose dispatch depends on inter-regional exchange
(Fig.~\ref{fig:robust}a). The first background field brings most of the amplification ($M^{(1)}/M^{(2)}$
between 0.73 and 0.97).

Threshold crossings, however, do not saturate as quickly, as expected for budget-dependent criticality: a small increase of $\Mik$ can move a field across $\tau$. At $\tau=0.7$ several
targets have fields that become critical only with a third background field: PJM net interchange (3
critical fields at $K=2$, 21 at $K=3$), NEM PPCCGT (10 and 24) and PJM load (14 and 22)
(Fig.~\ref{fig:robust}b). At $\tau=0.5$ all fields are already critical at $K=2$ except for TUMUT3 (22
versus 30). Taken together, these results show that small backgrounds recover most of the measured marginal risk, while the number of critical fields remains sensitive to the assumed budget. A stronger background model is evaluated by increasing $K$ and reporting the corresponding criticality results. Evaluating all sets of size four is
the workload for which the amortized scan is needed: for CAISO it comprises 52,360 sets per
target, against 7,175 for $K=2$.

\begin{figure}[!t]
\centering
\includegraphics[width=\textwidth]{fig9_robustness.pdf}
\caption{Robustness. (a) Saturation of mean field risk in the background budget. (b) Share of
$\tau$-critical fields at $\tau=0.7$ with $K=2$ (filled) and $K=3$ (open). (c) Agreement of reruns with
different attacker seeds or data splits with the main run. (d) Threshold sweep over ten targets:
fields that are safe alone but critical (solid), fields unsafe alone (dashed), and minimal unsafe sets
left after single-field grading (dotted).}
\label{fig:robust}
\end{figure}
